\documentclass[11pt,a4paper]{article}

\usepackage[utf8]{inputenc}
\usepackage[T1]{fontenc}
\usepackage{lmodern}
\usepackage[margin=0.76in]{geometry}
\usepackage{array}
\usepackage{tabularx}
\usepackage{booktabs}
\usepackage{enumitem}
\usepackage{titlesec}
\usepackage{fancyhdr}
\usepackage{lastpage}
\usepackage[hidelinks]{hyperref}

\hypersetup{
  pdftitle={Geography Standard Level Paper 1 Practice Simulation Set 2 Markscheme - 2026-05-07},
  pdfauthor={IB Geography SL Preparation}
}

\setlength{\parskip}{0.42em}
\setlength{\parindent}{0pt}
\setlength{\tabcolsep}{5pt}
\renewcommand{\arraystretch}{1.16}
\setlist[itemize]{leftmargin=1.32em,itemsep=3pt,topsep=4pt}

\newcolumntype{Y}{>{\raggedright\arraybackslash}X}
\newcolumntype{L}[1]{>{\raggedright\arraybackslash}p{#1}}

\titleformat{\section}
  {\normalfont\large\bfseries}
  {}{0pt}{}
\titleformat{\subsection}
  {\normalfont\normalsize\bfseries}
  {}{0pt}{}
\titleformat{\subsubsection}
  {\normalfont\normalsize\bfseries}
  {}{0pt}{}

\pagestyle{fancy}
\setlength{\headheight}{14pt}
\fancyhf{}
\fancyhead[C]{-- \thepage\ --}
\fancyhead[R]{Markscheme}
\renewcommand{\headrulewidth}{0pt}
\renewcommand{\footrulewidth}{0pt}

\newcommand{\term}[1]{\textbf{#1}}

\begin{document}

\thispagestyle{empty}
\begin{center}
{\Large Markscheme}\\[0.35em]
May 2026 practice\\[0.35em]
Geography\\[0.35em]
Standard level\\[0.35em]
Paper 1
\end{center}

\vspace{2em}

\section*{Paper 1 markbands}

These markbands are to be used for the 10-mark questions on paper 1 at both
standard level and higher level.

\begin{tabularx}{\textwidth}{L{0.09\textwidth}Y Y Y}
\toprule
\textbf{Marks}
& \textbf{AO1: Knowledge and understanding of specified content}\newline
\textbf{AO2: Application and analysis of knowledge and understanding}
& \textbf{AO3: Synthesis and evaluation}
& \textbf{AO4: Selection, use and application of a variety of appropriate skills and techniques} \\
\midrule
0 & \multicolumn{3}{Y}{The work does not reach a standard described by the descriptors below.} \\
1--2
& Response is very brief or descriptive, lists unconnected comments or largely
irrelevant information. Examples or case studies are not included or only
listed. Terminology is missing or used incorrectly.
& No evidence of evaluation or conclusion is expected at this level.
& Information is not grouped logically. Maps, graphs, or diagrams are not
included or are difficult to decipher, where appropriate. \\
3--4
& Response is general, with relevant and irrelevant information. Argument is not
relevant to the question. Basic terminology is used with errors.
& No evidence of critical evaluation. Conclusion, if appropriate, is irrelevant.
& Most information is not grouped logically. Maps, graphs, or diagrams lack
detail or are partially interpreted. \\
5--6
& Response describes relevant supporting evidence and outlines links to the
question. Argument partially addresses the question or repeats one point.
& Conclusion, if appropriate, is general or not aligned with the evidence. Other
perspectives or strengths and weaknesses of evidence are listed.
& Logically related information is grouped in sections or paragraphs but not
consistently. Maps, graphs, or diagrams do not follow conventions. \\
7--8
& Response describes relevant supporting evidence covering the main points of the
question. Argument is clear and relevant but one-sided or unbalanced. Complex
terminology is defined and used.
& Conclusion, if appropriate, is relevant and aligned with the evidence but
unbalanced. Other perspectives or strengths and weaknesses of evidence are
described.
& Logically related information is grouped in sections consistently. Maps,
graphs, or diagrams contribute to or support the argument. \\
9--10
& Response is in-depth and question-specific. Argument is balanced, presenting
evidence that is discussed, explaining complexity, exceptions, and comparisons.
Complex and relevant terminology is used correctly throughout.
& Conclusion, if appropriate, is relevant, balanced, and aligned with evidence.
Evaluation includes systematic presentation of ideas, cause and effect, other
perspectives, strengths and weaknesses, and justification of the argument.
& Response is logically structured with discussion focused on the argument or
points made. Maps, graphs, or diagrams are annotated following conventions and
their relevance is explained. \\
\bottomrule
\end{tabularx}

\newpage

\section*{Option F -- Food and health}

\subsection*{11.}

\subsubsection*{11. (a) (i) Identify the region with the highest percentage of food-insecure households. \hfill[1]}

\term{Dryland smallholder} or \term{dryland smallholder zone}. Allow \term{46\%}.

\subsubsection*{11. (a) (ii) State the percentage of households reporting child wasting in the dryland smallholder region. \hfill[1]}

\term{21\%}. Accept \term{20--22\%}.

\subsubsection*{11. (b) Outline one way in which an extended dry season can reduce safe water access in rural communities. \hfill[2]}

Award \term{[1]} for a valid way and \term{[1]} for development /
explanation.

For example, lower rainfall reduces recharge of shallow wells and surface
water sources \term{[1]}, leaving rural households reliant on more distant
or unsafe sources \term{[1]}.

Other possibilities include:
\begin{itemize}
  \item Falling river or reservoir levels raise the concentration of pathogens
  and pollutants in remaining water, lowering its safety.
  \item Damage to or congestion at functioning water points forces longer
  queues and unsafe storage, reducing effective access.
  \item Loss of livestock-watering points pushes households to share unsafe
  surface sources, raising contamination risk.
\end{itemize}

\subsubsection*{11. (c) (i) Explain how a climate shock can lead to changes in child wasting in vulnerable rural households. \hfill[3]}

Award \term{[1]} for a valid link and up to \term{[2]} for development /
explanation / exemplification.

For example, drought reduces crop yields and household food availability
\term{[1]}, lowering dietary diversity and energy intake for young children
\term{[1]}, leading to acute weight loss and child wasting in households
without food reserves \term{[1]}.

Other valid chains include:
\begin{itemize}
  \item Loss of livestock and farm income reduces money for high-quality foods
  and healthcare, and combined with diarrhoeal disease from unsafe water,
  worsens wasting in young children.
  \item Migration of working adults in search of income separates caregivers
  from children, reducing the frequency and quality of feeding and
  accelerating wasting in the youngest age groups.
\end{itemize}

\subsubsection*{11. (c) (ii) Explain how a climate shock can lead to changes in food prices in urban low-income districts. \hfill[3]}

Award \term{[1]} for a valid link and up to \term{[2]} for development /
explanation / exemplification.

For example, lower rural production cuts the supply of staple grains
\term{[1]}, raising wholesale and retail prices in urban markets \term{[1]},
which urban low-income households experience as reduced affordability and
poorer dietary quality \term{[1]}.

Other valid chains include:
\begin{itemize}
  \item Higher transport costs and panic buying after a shock further inflate
  urban food prices, forcing the poorest households to switch to cheaper, less
  nutritious foods.
  \item Currency depreciation or import-cost rises after a shock can compound
  domestic shortages, raising the price of imported staples in urban markets.
\end{itemize}

\newpage

\subsection*{12. (a) Examine the links between malnutrition and infectious disease. \hfill[10]}

Marks should be allocated according to the markbands.

Malnutrition, including undernutrition and micronutrient deficiency, weakens
immune function and increases the severity and duration of infectious disease.
Infectious disease, in turn, reduces appetite, nutrient absorption, and growth,
deepening malnutrition. Poverty, unsafe water, and conflict reinforce both
sides of this two-way link, producing a cycle that is difficult to break in
low-income or shock-affected communities.

Possible applied themes (AO2) demonstrating knowledge and understanding (AO1):
\begin{itemize}
  \item Undernutrition reduces immune competence, raising the incidence and
  severity of measles, pneumonia, diarrhoeal disease, and TB, for example
  child mortality in famine-affected Sahel or HIV-related wasting in southern
  Africa.
  \item Infectious disease, for example cholera in Haiti or intestinal
  parasites in West Africa, reduces nutrient absorption and increases nutrient
  loss, worsening malnutrition.
  \item Unsafe water and poor sanitation amplify both sides of the link by
  raising diarrhoeal disease, which is a leading cause of stunting and wasting
  in young children.
  \item Poverty, conflict, and displacement, for example Yemen or Tigray, act
  as common drivers, concentrating both undernutrition and infectious disease
  in the same populations.
  \item Public-health interventions such as vaccination, oral rehydration,
  WASH, school meals, and vitamin-A supplementation can break the cycle by
  cutting either side of the link.
\end{itemize}

Good answers may be well-structured (AO4) and may additionally offer a
critical evaluation (AO3) of the two-way relationship in named places,
recognising the role of poverty, water, sanitation, and health systems as
reinforcing factors, and may reach a balanced judgment.

For 5--6 marks, expect a weakly-evidenced outline of one link between
malnutrition and infectious disease, with at least one named example.

For 7--8 marks, expect a structured account which includes:
\begin{itemize}
  \item Either an evidenced examination of the two-way links between
  malnutrition and infectious disease in named places,
  \item Or a discursive conclusion, or ongoing evaluation, grounded in
  geographical concepts and/or perspectives.
\end{itemize}

For 9--10 marks, expect both of these traits.

\subsection*{12. (b) Evaluate the success of strategies used to improve food security following an environmental shock. \hfill[10]}

Marks should be allocated according to the markbands.

After an environmental shock, such as a drought, flood, or cyclone, food
security can be improved through short-term relief and longer-term resilience.
Relief includes food aid, cash transfers, and school meals. Longer-term
resilience includes irrigation, drought-resistant crops, storage, early-warning
systems, WASH, and microinsurance. Success varies with cost, equity, scale,
sustainability, and local fit, and depends on the speed and quality of
governance.

Possible applied themes (AO2) demonstrating knowledge and understanding (AO1):
\begin{itemize}
  \item Short-term relief, such as food aid, cash and voucher transfers, and
  school meals, can restore access quickly but may distort markets and create
  dependency if poorly designed.
  \item Production-side strategies, such as irrigation, drought-tolerant maize
  in sub-Saharan Africa, short-cycle crops, and crop diversification, raise
  resilience but require capital, water rights, and extension services.
  \item Storage and logistics, such as community grain banks, post-harvest
  loss reduction, and strategic reserves, stabilise supply but need governance
  and infrastructure.
  \item Safety nets, such as the Productive Safety Net Programme in Ethiopia,
  can support food security and human capital when transfers are linked to work
  or schooling.
  \item Early warning, WASH, and microinsurance address the shock-multiplier
  effects of disease, water, and income loss, but face uptake and equity
  challenges.
\end{itemize}

Good answers may be well-structured (AO4) and may additionally offer a
critical evaluation (AO3) of the success of strategies in named places,
recognising trade-offs between speed, cost, equity, and durability, and
reaching a supported overall judgment.

For 5--6 marks, expect a weakly-evidenced outline of one or more strategies
used after an environmental shock, with at least one named example.

For 7--8 marks, expect a structured account which includes:
\begin{itemize}
  \item Either an evidenced evaluation of the success of two or more
  food-security strategies after an environmental shock in named places,
  \item Or a discursive conclusion, or ongoing evaluation, grounded in
  geographical concepts and/or perspectives.
\end{itemize}

For 9--10 marks, expect both of these traits.

\newpage

\section*{Option G -- Urban environments}

\subsection*{13.}

\subsubsection*{13. (a) (i) State which zone has the highest share of trips made by walking or cycling. \hfill[1]}

\term{Informal hillside} or \term{informal hillside settlement}. Allow \term{45\%}.

\subsubsection*{13. (a) (ii) Identify which zone has the highest share of trips made by private car. \hfill[1]}

\term{Outer suburb}. Allow \term{74\%} or \term{outer car-dependent suburb}.

\subsubsection*{13. (b) Outline one reason why high private-car dependence in outer urban zones may create environmental stress. \hfill[2]}

Award \term{[1]} for a valid reason and \term{[1]} for development /
explanation.

For example, high private-car use raises tailpipe emissions of
NO\textsubscript{x} and PM\textsubscript{2.5} \term{[1]}, which lowers local
air quality and increases respiratory illness in residents \term{[1]}.

Other possibilities include:
\begin{itemize}
  \item Long suburban commutes increase fuel use and CO\textsubscript{2}
  emissions, contributing to the urban heat island and climate stress.
  \item Heavy traffic generates congestion and noise, lowering liveability and
  stressing road infrastructure.
  \item Car-oriented sprawl consumes farmland or habitat at the urban edge,
  reducing ecosystem services and water absorption.
\end{itemize}

\subsubsection*{13. (c) (i) Explain how the pattern of urban transport use can be affected by one economic factor. \hfill[3]}

Award \term{[1]} for an economic factor and up to \term{[2]} for development
/ explanation / exemplification.

For example, higher household incomes raise private-car ownership \term{[1]},
which encourages longer commutes and lower public-transport ridership
\term{[1]}, reinforcing car-based modal patterns even where transit is
available \term{[1]}.

Other economic factors include:
\begin{itemize}
  \item High fuel or parking costs reducing car use in favour of walking,
  cycling, or public transport.
  \item Fare subsidies or integrated ticketing lowering the cost of public
  transport and raising ridership.
  \item Employment location, for example a central business district, shaping
  the direction and length of trips.
\end{itemize}

\subsubsection*{13. (c) (ii) Explain how the pattern of urban transport use can be affected by one technological factor. \hfill[3]}

Award \term{[1]} for a technological factor and up to \term{[2]} for
development / explanation / exemplification.

For example, real-time mobile applications and contactless payments
\term{[1]} cut wait times and uncertainty for transit users \term{[1]},
encouraging public-transport use and reducing reliance on private cars
\term{[1]}.

Other technological factors include:
\begin{itemize}
  \item Bus rapid transit or metro investment, raising the passenger capacity
  of corridors and shifting modal split.
  \item E-bikes, e-scooters, and shared-mobility apps, extending door-to-door
  reach for short trips.
  \item Electric-vehicle charging networks shifting the environmental impact
  of car use rather than the modal split.
\end{itemize}

\newpage

\subsection*{14. (a) Examine the effectiveness of strategies used to manage urban traffic congestion. \hfill[10]}

Marks should be allocated according to the markbands.

Urban traffic congestion is the slowing of movement caused by demand exceeding
road capacity. It produces longer commutes, higher emissions, and higher
infrastructure costs. Strategies to manage congestion include demand management,
modal shift, and supply management. Their effectiveness varies with city form,
governance, equity, and time scale.

Possible applied themes (AO2) demonstrating knowledge and understanding (AO1):
\begin{itemize}
  \item Congestion is driven by rising vehicle ownership, dispersed origins and
  destinations, peak-hour demand, and limited road space.
  \item Demand-management tools, such as London congestion charging or
  Singapore Electronic Road Pricing, reduce vehicle entry and raise revenue,
  but raise equity concerns for lower-income drivers.
  \item Modal-shift strategies, such as Bogot\'{a} or Curitiba bus rapid
  transit, Tokyo rail, or Paris cycling networks, move passengers per hour at
  lower cost than private cars, but require dense networks and steady
  investment.
  \item Supply-management strategies, such as road expansion, ring roads, and
  intelligent traffic systems, can release short-term capacity but often induce
  demand and worsen sprawl.
  \item Effectiveness depends on integration with land-use planning, political
  commitment, public acceptance, and durability over time.
\end{itemize}

Good answers may be well-structured (AO4) and may additionally offer a
critical evaluation (AO3) of the relative effectiveness of demand-side and
supply-side strategies, in different cities and at different time scales,
recognising trade-offs between cost, equity, and environmental outcome.

For 5--6 marks, expect a weakly-evidenced outline of one or more
congestion-management strategies in at least one named city.

For 7--8 marks, expect a structured account which includes:
\begin{itemize}
  \item Either an evidenced examination of the effectiveness of two or more
  congestion-management strategies in named cities,
  \item Or a discursive conclusion, or ongoing evaluation, grounded in
  geographical concepts and/or perspectives.
\end{itemize}

For 9--10 marks, expect both of these traits.

\subsection*{14. (b) To what extent does urban sprawl create greater environmental challenges than high-density urban growth? \hfill[10]}

Marks should be allocated according to the markbands.

Urban sprawl is the low-density, often discontinuous expansion of urban land at
the periphery, typically dependent on private cars. High-density urban growth
concentrates population and activity in compact built environments. Both create
environmental challenges, but of different kinds and at different scales. The
extent of greater challenge depends on how environmental challenge is defined,
including air pollution, land use, heat, water, biodiversity, and emissions per
capita.

Possible applied themes (AO2) demonstrating knowledge and understanding (AO1):
\begin{itemize}
  \item Sprawl is associated with car dependence, higher per-capita emissions,
  loss of farmland and habitat, longer water and energy networks, and
  segregation by income, for example in U.S. Sun Belt cities or peri-urban
  Lagos.
  \item High-density growth is associated with congestion, localised air
  pollution and noise, urban heat island, overcrowding, pressure on green
  space, and high land values, for example Hong Kong or central Mumbai.
  \item Per-capita emissions and land-use efficiency often favour high-density
  growth, while local liveability and air quality often favour lower-density
  growth, complicating the comparison.
  \item Time scale matters: sprawl locks in long-term car infrastructure;
  density locks in vertical infrastructure that is also hard to retrofit.
  \item Governance, planning, and investment in transit, green space, and
  building standards can change the size of either set of challenges.
\end{itemize}

Good answers may be well-structured (AO4) and may additionally offer a
critical evaluation (AO3) of which form creates greater challenges, using named
examples and explicit criteria, and may reach a balanced judgment that
recognises trade-offs and the role of planning.

For 5--6 marks, expect a weakly-evidenced outline of challenges from sprawl
and/or high-density growth in one or more named places.

For 7--8 marks, expect a structured account which includes:
\begin{itemize}
  \item Either an evidenced examination of challenges from both sprawl and
  high-density growth in named places,
  \item Or a discursive conclusion, or ongoing evaluation, grounded in
  geographical concepts and/or perspectives.
\end{itemize}

For 9--10 marks, expect both of these traits.

\end{document}
